\subsection{\identifier{UC-13} --- Upload and Update Worker Documents}
\label{uc:13}

\paragraph{Purpose}

Allow an authorised user to upload documents
and create new versions associated with
one worker.

This use case covers existing files supplied for upload. Generation follows
\identifier{UC-20} and signing follows \identifier{UC-14}; the workflows share the
same version and retention rules.


\paragraph{Actors}

\textbf{Primary actor:} Farm Owner.

\textbf{Additional actor:} Field Manager,
with upload or update
permissions for the relevant document type.

Insurance consultation permission does not
automatically permit document modification.

\paragraph{Preconditions}

\begin{itemize}
    \item An internet connection is available.
    \item The user has permission
    for the document type and affected workers.
    \item The selected worker record exists within the authorised business.
    \item For a version update,
    the target document exists.
    \item A supported file is available.
\end{itemize}

\paragraph{Main Flow}

\begin{enumerate}
    \item The user opens document management.
    \item The user chooses a new document
    or a new version of an existing document.
    \item The user selects or reviews the document's single worker.
    \item The user selects the document category:
    contract, insurance or other, and records a form subtype where known
    (Table~\ref{tab:source-forms}).
    \item The user provides or confirms
    the document title and selects the file.
    \item For insurance, the user enters
    the expiry date where known.
    \item The user may add observations
    and declare whether the uploaded version
    was signed outside the application.
    \item The system validates permissions,
    required fields, file format and size.
    \item The system presents the submission
    and relevant version-retention consequences.
    \item The user confirms.
    \item The system stores the accepted version,
    records its upload date and version identity,
    and makes it current.
    \item The system applies the retention policy
    and confirms the outcome.
\end{enumerate}

\paragraph{Alternative Flows and Exceptions}

\begin{description}

    \item[A1 --- An insurance expiry date is unknown.]
    The system allows the document to be saved
    with validity identified as unconfirmed.
    It does not infer validity from upload date.

    \item[A2 --- The selected worker is incorrect.]
    The user corrects the selection before acceptance. An accepted document's
    worker identity is fixed; a document for another worker is a separate upload.

    \item[A3 --- A version was signed externally.]
    The user explicitly identifies that status
    for the uploaded version.
    It is distinguishable from an application
    signing event and is not presented
    as externally verified by the system.

    \item[A4 --- An existing signed document is updated.]
    The new version becomes current
    without inheriting an earlier signature.
    Successful replacement deletes the old source and signing files
    under A5.


    \item[A5 --- A replacement is accepted.]
    For every category, including signed contracts, keep the current source until
    the new upload is complete and validated. Publish the replacement and delete
    the previous source, signature, signed PDF and external-signature declaration.
    The replacement is unsigned unless explicitly declared externally signed in
    its own upload; it inherits no status. Keep no consultable version history.
    Pending signing of the superseded source requires rejection/recovery under
    \decisionref{Q-DOC-04}; it cannot delay replacement indefinitely or recreate
    deleted content. Propagate deletion to prepared copies on synchronisation.

    \item[A6 --- An incorrect upload may be deleted.]
    An authorised user may confirm deletion
    when the target has no signature
    and no dependent records.

    \item[A7 --- An incorrect upload is signed.]
    The user may withdraw it from use,
    preserving its file and signature association
    while it remains retained.
    The current withdrawn version remains until successful replacement, which deletes its file, signature and signed PDF under \identifier{BR-026}.

    If the withdrawn version is current, it remains identified
    as current and withdrawn until a replacement is accepted.

    An earlier version does not automatically become current.

    \item[A8 --- Validation fails.]
    The system identifies unsupported formats,
    files exceeding the size limit,
    missing required information
    or invalid associations.
    No version is accepted until corrected.

    \item[A9 --- Upload or storage fails.]
    The system reports the failure.
    The current source and its signing files remain unchanged.
    Failed uploads do not trigger retention deletion.

    \item[A10 --- A submission is repeated.]
    Retrying the same upload operation
    does not unintentionally publish
    duplicate versions or trigger extra deletion.

    \item[A11 --- Connectivity is unavailable.]
    Upload and update operations cannot complete.
    Prepared offline insurance consultation
    remains available through \identifier{UC-12}.

\end{description}

\paragraph{Postconditions}

\textbf{On success:}
The accepted version becomes current
and is available for authorised consultation
through the associated worker records.


\textbf{On any successful replacement:}
Only the new current source remains available for the document. The old source,
signature and signed PDF are deleted, including for contracts; no signing status
is inherited and no version-history browser is provided.

\textbf{On failure or cancellation:}
The previously accepted current source and its signing files are unchanged
and no retention deletion is performed.

\textbf{In all cases:}
A new version does not inherit signatures
from another version.

\paragraph{Related Requirements and Rules}
\hyperref[fr:040]{\identifier{FR-040}}, \hyperref[fr:041]{\identifier{FR-041}}, \hyperref[fr:044]{\identifier{FR-044}}, \hyperref[fr:063]{\identifier{FR-063}}.

\hyperref[br:022]{\identifier{BR-022}}, \hyperref[br:023]{\identifier{BR-023}}, \hyperref[br:024]{\identifier{BR-024}}, \hyperref[br:025]{\identifier{BR-025}}, \hyperref[br:026]{\identifier{BR-026}}, \hyperref[br:027]{\identifier{BR-027}}, \hyperref[br:028]{\identifier{BR-028}}, \hyperref[br:030]{\identifier{BR-030}}.

\paragraph{Decision References}
\decisionref{Q-DOC-06}, \decisionref{D-07}, \decisionref{Q-DOC-07}, \decisionref{Q-DOC-01}, \decisionref{Q-DOC-04}.

\paragraph{Uploaded examples}
The four individual forms may be uploaded here. Retain the
original file. An external-signature declaration does not imply acceptance of a
medical examination. A structured response may be recorded only when explicitly
visible in the source and reviewed; ambiguous responses remain unknown.
